\section{Literature review: evidence framework}

This section is a framework to be populated \emph{after} collecting and verifying
primary sources. The companion \texttt{sources/evidence\_matrix.csv} records each
source's link, modality, sample size, results, limitations, and verification level.

\subsection{Conformal-coating process and acceptability}
Review coating chemistries, fluorescent additives, defect mechanisms and applicable
IPC standards. Separate an optical indication of coating from an independently
validated thickness or functional reliability measurement.

\subsection{Fluorescence optics and measurement}
Examine excitation/emission spectra, UV blocking versus visible collection,
illumination geometry, exposure linearity, shadows and flat-field calibration.

\subsection{Classical inspection and geometric reference models}
Compare image registration, thresholding, color segmentation, morphology,
connected components and CAD/golden-board required/forbidden masks.

\subsection{Learned segmentation and anomaly detection}
U-Net introduced an encoder--decoder segmentation architecture
\citep{ronneberger2015unet}. SegFormer is a transformer-based semantic
segmentation architecture \citep{xie2021segformer}. PatchCore is a general
industrial anomaly-detection approach \citep{roth2022patchcore}.
These are \emph{method baselines}; the cited papers are not UV PCB coating
validation studies.

\subsection{Datasets, reproducibility and edge deployment}
Search for accessible UV-coated-PCB data separately from conventional bare-PCB,
solder or assembly-defect benchmarks. Review specimen-level data leakage,
external validation, CPU deployment and end-to-end inspection timing.

\subsection{Evidence gap ledger}
Do not assert a novelty claim until a documented search of direct conformal
coating, fluorescence imaging, industrial AOI and adjacent methods is complete.
